\section{What this says about agent systems}
\label{sec:agents}

\subsection{The free compression ratio}

Define the \emph{compression ratio} of a flat record,
$C(n) = 2^n / \Kflat(n)$: how much smaller the record is than the full effort
profile. Dividing the bounds of Section~\ref{sec:threshold} through by $2^n$,
\[
  2^{\lfloor \log_2(n/2+1)\rfloor} \;\le\; C(n) \;\le\; \frac{n+2}{2}.
\]
Both ends are $\Theta(n)$, so the safe zone does widen as the team grows---but
only linearly, against an exponentially large baseline.

\begin{table}[t]
\centering
\caption{The same numbers under two normalisations. $(n+2)/2$ is an upper bound
on $C(n)$, not its value; the measured ratio sits at or below it, and the gap is
real at $n=7$.}
\label{tab:ratio}
\begin{tabular}{@{}lrrrrrr@{}}
\toprule
$n$ & 3 & 5 & 7 & 9 & 20 & 50 \\
\midrule
$C(n)$        & 2.00 & 3.20 & 4.00 & 4.00--5.28 & 8--11 & 16--26 \\
$(n+2)/2$     & 2.50 & 3.50 & 4.50 & 5.50 & 11 & 26 \\
free bits     & 1.00 & 1.68 & 2.00 & 2.00--2.40 & 3.00--3.46 & 4.00--4.70 \\
as \% of $n$  & 33.3 & 33.6 & 28.6 & 22--27 & 15.0--17.3 & 8.0--9.4 \\
\bottomrule
\end{tabular}
\end{table}

\subsection{Two normalisations, both correct}

Measured as a \emph{state-space ratio}, the safe zone widens: at $n=5$ one may
compress $3.2\times$ for free, at $n=50$ about $26\times$. Measured in
\emph{bits}---the unit a record is actually sized in---the free allowance is
$\log_2 C(n) \approx \log_2(n/2)$ bits out of $n$, so the \emph{fraction} of
per-agent resolution one may discard for free \emph{shrinks} with team size
(Table~\ref{tab:ratio}). Linear growth of a ratio against an exponential baseline
is logarithmic growth in bits; both statements describe the same numbers.

Everything in this section is computed under the uniform prior. What a
non-uniform prior does to these numbers we do not know:
Section~\ref{sec:priors} shows that a tilted prior can move the threshold in
either direction, and leaves the general case open. Read the table as the
uniform-prior answer, not as a bound on every prior.

Concretely: at $n = 20$, flatness requires $K$ between $95{,}326$ and $131{,}072$
labels---retaining $16.5$--$17.0$ of the $20$ bits, i.e.\ $83$--$85\%$ of
per-agent resolution. At $n = 50$ it requires retaining $45.3$--$46.0$ of the
$50$ bits, i.e.\ $91$--$92\%$.

\subsection{Consequences for system design}

\paragraph{Keep one bit per agent and the problem disappears.}
A record that retains a single effort bit per agent is flat, with margin
$1$, and no tension exists. This is the first thing the theory says, and it is
worth saying plainly: the results below are about what happens when that is not
affordable.

\paragraph{Aggregates are far inside the tension regime.}
A statistic such as ``how many agents succeeded'' carries $\log_2(n+1)$ bits---
$5.7$ bits at $n = 50$, about $11\%$ of per-agent resolution---which is an order
of magnitude below the flatness threshold. By Theorem~\ref{thm:elim}\,(ii), any
record that maps a profile and all $n$ of its unilateral neighbours to one label
admits no positive margin under any bounded payment.

\paragraph{Summaries that aggregate across agents.}
The compaction step described in Section~\ref{sec:intro} replaces a $k$-agent
fan-out with a paragraph. We have not measured how many bits of per-agent
resolution such a paragraph retains, and the theory cannot tell us: what it gives
is the bar such a record has to clear. A $20$-agent episode escapes the tension
only by retaining $16.5$ of $20$ bits of \emph{per-agent} resolution under the
most favourable prior, and essentially all $20$ under any other. Whenever a
summarisation step falls short of that---a question about a particular system,
settled by measuring it---the two uses of the record genuinely trade off, and
Theorem~\ref{thm:onebit} gives the rate of that trade at the extreme end.
Section~\ref{sec:experiment} measures where a family of concrete compaction
mechanisms lands.

\paragraph{Where the trade is made deliberately.}
Bittensor's concealment interval~\cite{bittensorcommitreveal}, tuned to roughly $1800$ blocks, is the clearest
deployed instance of an operator selecting a point on this curve: more
concealment, more incentive margin, slower evaluation. We note honestly that its
binding constraint is temporal and cryptoeconomic rather than a size bound---it
evidences the \emph{consequence} of settling from a merged record, while the
context-window argument above evidences the \emph{mechanism}. Neither alone
makes the case; together they do.

\subsection{What production agent frameworks log by default}
\label{sec:framework-audit}

The theory bounds what any record must retain; whether real systems clear that
bar is a separate, checkable question. We audited seven production
multi-agent frameworks against one criterion: does the framework's
\emph{default} structured log let a principal recover, after the fact, which
agent did what---not whether the information is technically reconstructable
with custom instrumentation, but what a settlement built on the log as shipped
would see. For AutoGen we went a step further than reading documentation: we
installed \texttt{autogen-core}~0.7.5 and instrumented a real three-agent
episode through its own \texttt{SingleThreadedAgentRuntime}, capturing every
event the framework's own logger actually emitted, and checked the live
objects rather than the API reference page (\texttt{experiments/production\_probe/}
in the artifact reproduces the run).

\begin{table}[t]
\centering
\caption{Default per-agent attribution in seven production agent frameworks,
audited 2026-09-14, re-verified 2026-09-19, and for AutoGen corrected
2026-09-19 after running the framework directly (see text).}
\label{tab:frameworks}
\begin{tabular}{@{}lll@{}}
\toprule
Framework & Class & Effort recoverable? \\
\midrule
LangSmith / LangGraph~\cite{langsmithdocs} & SAFE & per-step, per-step cost \\
OpenAI Agents SDK~\cite{openaiagentssdk} & SAFE & per agent run \\
Claude Agent SDK & SAFE & per subagent, own file \\
Microsoft Semantic Kernel~\cite{semantickerneltokens} & SAFE & per-agent token/cost \\
Google Agent Development Kit~\cite{googleadktraces} & SAFE (leaning) & per-agent spans \\
AutoGen~\cite{autogenlogging} & SAFE & per agent, via ambient context \\
CrewAI~\cite{crewaitracing,crewaiovercountissue} & MIXED & outputs yes, cost no \\
\bottomrule
\end{tabular}
\end{table}

Six of the seven default to per-agent attribution; only one is mixed. This is
a correction, not a stylistic update: an earlier version of this audit, based
on the AutoGen API reference page, classified AutoGen MIXED on the strength of
\texttt{LLMCallEvent} and \texttt{ToolCallEvent} appearing to carry no agent
identifier in their public constructor signature. Running the framework
directly shows this was wrong. Both classes populate an \texttt{agent\_id}
field internally from \texttt{MessageHandlerContext.agent\_id()} whenever the
event is logged from inside a message handler---the normal case for a real
agent's own LLM or tool call---and this field is undocumented on the public
reference page, which is why the documentation-only audit missed it. In our
instrumented run, 120 of 120 \texttt{LLMCallEvent}s and 120 of 120
\texttt{ToolCallEvent}s across 40 episodes carried a working per-agent
identifier; \texttt{MessageEvent}'s sender/receiver fields, the basis of the
original SAFE-adjacent reading, are unaffected. We leave this correction
visible rather than quietly fixing the table, because it is itself evidence
for a point the paper should make honestly: \emph{auditing framework logging
behavior from documentation alone is not reliable}, which is a reason for
skepticism about any claim in this section, including our own, that has not
been checked against a live run. CrewAI remains the one mixed case we found:
it keeps per-task outputs but exposes token usage crew-wide; when multiple
agents share one LLM client, which is the documented usual way to write a
crew, the aggregate is not merely uninformative but actively overcounted
relative to any individual agent's
share~\cite{crewaiovercountissue}.

\paragraph{This cuts against the easy version of our own motivation.} If the
industry default were aggregation, the case for this paper would be
academic. It is not: most production frameworks keep per-agent resolution
without being asked to. The honest reading is narrower than that would
suggest, and we do not soften it. In every mixed case found, the aggregation
is a convenience default, not a binding budget---the missing attribution is
recoverable with modest instrumentation, and in CrewAI's case practitioners
already build third-party tooling to retrofit it, which is itself evidence
that the gap is felt. That is a weaker premise than the genuinely bounded
record this paper's theory targets, and Section~\ref{sec:limitations} says so
without qualification. What the audit establishes is narrower and, we think,
still useful: the failure mode this paper formalises is not hypothetical---it
already exists at the schema level in two widely used frameworks, on exactly
the axis the theory predicts, and it is currently invisible to anyone who has
not gone looking at the field level.
